% Luc Destombe --- licence CC BY-NC-SA 4.0
% https://creativecommons.org/licenses/by-nc-sa/4.0/deed.fr
% Fichier autonome : compiler avec pdflatex, deux fois (signets, nombre de pages).
\documentclass[11pt,a4paper]{article}
\makeatletter
\newif\ifeleve
\elevefalse

\newif\iflycee@palatino
\lycee@palatinotrue

\RequirePackage[utf8]{inputenc}
\RequirePackage[T1]{fontenc}
\RequirePackage[french]{babel}
\RequirePackage{etoolbox}

\newcommand{\ShorthandsOff}{\@ifundefined{shorthandoff}{}{\shorthandoff{;:!?}}}
\newcommand{\ShorthandsOn}{\@ifundefined{shorthandon}{}{\shorthandon{;:!?}}}
\ShorthandsOff
\AtBeginDocument{\ShorthandsOn}
\AtBeginEnvironment{tikzpicture}{\ShorthandsOff}

\RequirePackage{geometry}
\RequirePackage{fancyhdr}
\RequirePackage{lastpage}
\RequirePackage{multicol}
\RequirePackage{array}
\RequirePackage{enumitem}
\RequirePackage{xcolor}
\RequirePackage{needspace}

\RequirePackage{amsmath,amssymb}
\RequirePackage{mathpazo}
\DeclareMathSizes{10.95}{10.95}{8}{5.5}
\begingroup\catcode`\_=\active \catcode`\^=\active
\gdef_#1{\sb{\scriptscriptstyle #1}}
\gdef^#1{\sp{\scriptscriptstyle\lycee@moinsepais #1}}
\endgroup
\newcommand\lycee@moins{\mathbin{%
  \setbox\z@\hbox{$\m@th\scriptscriptstyle\lycee@moinsorig$}%
  \dimen@=.55\wd\z@ \advance\dimen@-.5pt
  \hbox to.75\wd\z@{\hss$\m@th\scriptscriptstyle\vcenter{\hbox{%
    \tikz\draw[line width=.5pt,line cap=round](0,0)--(\the\dimen@,0);}}$\hss}}}
\begingroup\lccode`\~=`\- \lowercase{\endgroup\let~}\lycee@moins
\newcommand\lycee@moinsepais{\mathcode`\-="8000 }
\AtBeginDocument{\mathchardef\lycee@moinsorig=\mathcode`\-\relax}
\AtBeginDocument{\catcode`\_=12 \mathcode`\_="8000
  \catcode`\^=12 \mathcode`\^="8000 }
\newcommand\lycee@exposants{%
  \fontdimen13\textfont2=.48\fontdimen6\textfont2
  \fontdimen14\textfont2=.48\fontdimen6\textfont2
  \fontdimen15\textfont2=.48\fontdimen6\textfont2
  \fontdimen13\scriptfont2=.48\fontdimen6\scriptfont2
  \fontdimen14\scriptfont2=.48\fontdimen6\scriptfont2
  \fontdimen15\scriptfont2=.48\fontdimen6\scriptfont2}
\every@math@size\expandafter{\the\every@math@size\lycee@exposants}
\newlength{\retraitcalculs}
\setlength{\retraitcalculs}{2em}

\RequirePackage{tikz}
\usetikzlibrary{arrows.meta,calc,babel,shapes.misc}
\RequirePackage[most]{tcolorbox}
\tcbuselibrary{skins,breakable}

\definecolor{coulDef}{HTML}{1F4E79}
\definecolor{coulProp}{HTML}{7030A0}
\definecolor{coulRem}{HTML}{C55A11}
\definecolor{coulEx}{HTML}{2E7D32}
\definecolor{coulExo}{HTML}{B03A2E}
\definecolor{coulPart}{HTML}{1F4E79}
\definecolor{coulMeth}{HTML}{1B6E4F}
\definecolor{coulCourbe}{HTML}{0B5ED7}
\definecolor{coulCourbeB}{HTML}{C00000}
\definecolor{coulCourbeC}{HTML}{2E7D32}
\definecolor{coulGrille}{HTML}{8C8C8C}
\definecolor{coulRep}{HTML}{1B6E4F}
\definecolor{coulBar}{HTML}{2E7D32}

\newcommand{\R}{\mathbb{R}}
\newcommand{\Rs}{\mathbb{R}^{*}}
\renewcommand{\le}{\leqslant}
\renewcommand{\ge}{\geqslant}
\renewcommand{\approx}{\simeq}
\newcommand{\vect}[1]{\overrightarrow{#1}}
\newcommand{\norme}[1]{\left\lVert #1 \right\rVert}
\newcommand{\intervalle}[2]{\left[#1\,;#2\right]}
\RequirePackage{cancel}
\renewcommand{\CancelColor}{\color{coulExo}}
\newcommand{\fois}[1]{\times{\color{coulExo}#1}}
\newcommand{\barre}[1]{\cancel{#1}}

\tikzset{grille/.style={coulGrille,line width=0.4pt}}
\newcommand{\quadrillage}[4]{\draw[grille,step=1] (#1,#2) grid (#3,#4);}
\tikzset{
  croix/.style={cross out,draw,line width=0.6pt,inner sep=0pt,minimum size=4pt},
  croixdroite/.style={croix,rotate=45,line width=0.8pt},
}
\newcommand{\pt}[3]{\path (#1) node[croix]{} node[#2,font=\small,inner sep=2.5pt]{$#3$};}

\newcommand{\lycee@repere}[4]{%
  \draw[grille,step=1] (#1,#2) grid (#3,#4);
  \draw[-{Stealth[length=2mm]},thick] (#1,0)--({#3+0.4},0) node[below left=-1pt]{$x$};
  \draw[-{Stealth[length=2mm]},thick] (0,#2)--(0,{#4+0.4}) node[below left=-1pt]{$y$};
  \node[below left,font=\scriptsize,inner sep=1.5pt] at (0,0) {$0$};}
\newcommand{\lycee@gradx}[1]{\foreach \k/\l in {#1}{%
  \draw (\k,0.07)--(\k,-0.07) node[below,font=\scriptsize,inner sep=1.5pt]{$\l$};}}
\newcommand{\lycee@grady}[1]{\foreach \k/\l in {#1}{%
  \draw (0.07,\k)--(-0.07,\k) node[left,font=\scriptsize,inner sep=1.5pt]{$\l$};}}
\newcommand{\lycee@intff}[2]{\left[#1\,;#2\right]}
\newcommand{\lycee@intfo}[2]{\left[#1\,;#2\right[}
\newcommand{\lycee@intof}[2]{\left]#1\,;#2\right]}
\newcommand{\lycee@intoo}[2]{\left]#1\,;#2\right[}
\newcommand{\lycee@ens}[1]{\left\{#1\right\}}
\AtBeginDocument{%
  \providecommand{\repere}{\lycee@repere}%
  \providecommand{\gradx}{\lycee@gradx}%
  \providecommand{\grady}{\lycee@grady}%
  \providecommand{\Cf}{\mathcal{C}\sb{\scriptscriptstyle f}}%
  \providecommand{\Cg}{\mathcal{C}\sb{\scriptscriptstyle g}}%
  \providecommand{\intff}{\lycee@intff}%
  \providecommand{\intfo}{\lycee@intfo}%
  \providecommand{\intof}{\lycee@intof}%
  \providecommand{\intoo}{\lycee@intoo}%
  \providecommand{\ens}{\lycee@ens}}

\newlist{questions}{enumerate}{3}
\setlist[questions,1]{label=\arabic*\textdegree),leftmargin=*,itemsep=2pt,topsep=3pt}
\setlist[questions,2]{label=\alph*),leftmargin=*,itemsep=1pt,topsep=2pt}
\setlist[questions,3]{label=\roman*),leftmargin=*,itemsep=1pt,topsep=1pt}
\setlist[itemize]{label=\textbullet,leftmargin=*,itemsep=2pt,topsep=3pt}

\newcommand{\besoinplace}[1]{\par\penalty\@M\begingroup
  \dimen@=#1\relax\dimen@ii=\pagegoal\advance\dimen@ii-\pagetotal
  \ifdim\dimen@>\dimen@ii\ifdim\dimen@ii>\z@\vfil\fi\break\fi\endgroup}

\newcommand{\lycee@pied}{\thepage/\pageref{LastPage}}
\newcommand{\entete}[2]{%
  \pagestyle{fancy}\fancyhf{}%
  \fancyhead[L]{\footnotesize\color{coulPart}\textbf{#1}}%
  \fancyhead[R]{\footnotesize\color{coulPart}\textbf{#2}}%
  \fancyfoot[C]{\footnotesize\lycee@pied}%
  \renewcommand{\headrulewidth}{0.4pt}%
  \renewcommand{\headrule}{\hbox to\headwidth{%
    \color{coulPart!40}\leaders\hrule height \headrulewidth\hfill}}}

\RequirePackage{ccicons}
\newcommand{\lycee@licence}{{\scriptsize\color{gray}%
  \copyright\ Luc Destombe --- \ccbyncsa\ CC BY-NC-SA 4.0}}
\newif\iflycee@licence \lycee@licencetrue
\newcommand{\sanslicence}{\lycee@licencefalse}
\AtBeginDocument{\iflycee@licence\AddToHookNext{shipout/foreground}{%
  \put(0,0){\rlap{\kern\dimexpr1in+\hoffset+\oddsidemargin+\textwidth\relax
    \llap{\raisebox{-\dimexpr1in+\voffset+\topmargin+\headheight+\headsep
      +\textheight+\footskip\relax}{\lycee@licence}}}}}\fi}

\newcommand{\formulesserrees}{%
  \setlength{\abovedisplayskip}{3pt}\setlength{\belowdisplayskip}{3pt}%
  \setlength{\abovedisplayshortskip}{2pt}\setlength{\belowdisplayshortskip}{3pt}}

\setlength{\parindent}{0pt}

\RequirePackage[implicit=false,hidelinks,pdfencoding=auto,psdextra,bookmarksopen,
  bookmarksopenlevel=1,pdfcreator={},pdfproducer={}]{hyperref}
\RequirePackage{bookmark}
\hypersetup{pdfauthor={Luc Destombe},pdfkeywords={CC BY-NC-SA 4.0}}
\newcounter{lycee@signet}
\newcommand{\signet}[3][]{\stepcounter{lycee@signet}%
  \edef\lycee@dest{\if\relax\detokenize{#1}\relax
    signet.\arabic{lycee@signet}\else#1\fi}%
  \hypertarget{\lycee@dest}{}\bookmark[level=#2,dest=\lycee@dest]{#3}}
\pdfstringdefDisableCommands{%
  \def\R{R}\def\vect#1{#1}\def\Cf{Cf}\def\Cg{Cg}\def\fois#1{×#1}%
  \def\barre#1{#1}\def\intervalle#1#2{[#1;#2]}\def\intff#1#2{[#1;#2]}%
  \def\textsuperscript#1{#1}\def\dfrac#1#2{#1/#2}\def\frac#1#2{#1/#2}%
  \def\vec#1{#1⃗}}%
\geometry{top=2cm,bottom=1cm,left=1cm,right=1cm,headheight=14pt,footskip=0.6cm}

\RequirePackage{pgfplots}
\pgfplotsset{compat=1.18}
\RequirePackage{tkz-tab}

\newcounter{partiecours}
\renewcommand{\thepartiecours}{\Roman{partiecours}}
\newcommand{\partiecours}[1]{%
  \refstepcounter{partiecours}%
  \Needspace*{8\baselineskip}%
  \bigskip
  \noindent\begin{tcolorbox}[enhanced,colback=coulPart,colframe=coulPart,
      boxrule=0pt,arc=2pt,left=8pt,right=8pt,top=4pt,bottom=4pt,
      before skip=6pt,after skip=10pt]
    \color{white}\bfseries\large \thepartiecours{} --- #1%
    \signet[partie-\arabic{partiecours}]{1}{\thepartiecours{} --- #1}
  \end{tcolorbox}%
  \addcontentsline{toc}{section}{\thepartiecours{} --- #1}%
}

\newtcolorbox{boitecours}[3][]{%
  enhanced,breakable,
  colback=white,colframe=#2,coltitle=white,
  fonttitle=\bfseries,
  attach boxed title to top left={xshift=8pt,yshift=-\tcboxedtitleheight/2},
  boxed title style={colback=#2,arc=2pt,boxrule=0pt},
  boxrule=0.9pt,arc=2pt,
  left=8pt,right=8pt,top=8pt,bottom=6pt,
  title={#3},#1}

\newcounter{methode}
\newenvironment{definitionbox}[1][Définition]%
  {\begin{boitecours}[phantom={\signet{2}{#1}}]{coulDef}{#1}}{\end{boitecours}}
\newenvironment{proprietebox}[1][Propriété]%
  {\begin{boitecours}[phantom={\signet{2}{#1}}]{coulProp}{#1}}{\end{boitecours}}
\newenvironment{methodebox}[1][Méthode]{\stepcounter{methode}%
  \begin{boitecours}[phantom={\signet[methode-\arabic{methode}]{2}{#1}}]{coulMeth}{#1}}%
  {\end{boitecours}}

\newcommand{\etape}[1]{\par\smallskip\textbf{\color{coulMeth}#1}\par\nobreak\smallskip}
\newcommand{\enonce}[1]{\emph{#1}\par\smallskip}
\newcommand{\reponse}[1]{\par\smallskip
  {\footnotesize\itshape\color{black!55}Réponses : #1}}
\newcommand{\demo}[1]{\textbf{\color{coulMeth}Démonstration.} #1}
\newcommand{\afaire}[1]{%
  \par\smallskip
  \noindent{\small\color{coulExo}$\blacktriangleright$\ \textbf{À faire :}
    fiche d'exercices du chapitre, #1.}\par\smallskip}

\newenvironment{calculs}{\setlength{\jot}{5pt}\@fleqntrue\@mathmargin\retraitcalculs
  \setlength{\abovedisplayskip}{4pt}\setlength{\belowdisplayskip}{4pt}%
  \setlength{\abovedisplayshortskip}{4pt}\setlength{\belowdisplayshortskip}{4pt}%
  \csname alignat*\endcsname{2}}%
  {\csname endalignat*\endcsname}
\newcommand{\rem}[1]{\text{\small\itshape\color{black!60}#1}}
\newcommand{\explique}[2]{\par\smallskip\noindent
  \begin{minipage}[t]{0.57\linewidth}\raggedright #1\end{minipage}\hfill
  \begin{minipage}[t]{0.40\linewidth}\raggedright\leavevmode
    {\small\itshape\color{black!60}#2}\end{minipage}\par}
\newcommand{\cas}[1]{\par\smallskip\noindent\textbf{\color{coulExo}#1}\nobreak}
\newcommand{\calc}[1]{\par\smallskip\textbf{\color{coulExo}#1}\par\nobreak}

\newtcolorbox{remarquebox}[1][Remarque]{%
  enhanced,breakable,colback=white,colframe=coulRem,
  boxrule=0pt,leftrule=2.5pt,arc=0pt,outer arc=0pt,
  left=8pt,right=6pt,top=5pt,bottom=5pt,
  before upper={\textbf{\color{coulRem}#1}\par\smallskip}}

\newtcolorbox{exemplebox}[1][Exemple]{%
  enhanced,breakable,colback=white,colframe=coulEx,
  boxrule=0pt,leftrule=2.5pt,arc=0pt,outer arc=0pt,
  left=8pt,right=6pt,top=5pt,bottom=5pt,
  before upper={\textbf{\color{coulEx}#1}\par\smallskip}}

\pgfplotsset{
  repere/.style={
    axis lines=middle,
    axis line style={-{Stealth[length=2.4mm]},thick},
    every axis x label/.style={at={(ticklabel* cs:1.0)},anchor=west},
    every axis y label/.style={at={(ticklabel* cs:1.0)},anchor=south},
    label style={font=\footnotesize},
    tick label style={font=\scriptsize},
    grid=both,
    major grid style={grille},
    minor tick num=0,
    tick align=inside,
    clip mode=individual,
  },
}
\tikzset{
  courbe/.style={coulCourbe,line width=1.1pt,smooth},
  courbeB/.style={coulCourbeB,line width=1.1pt,smooth},
  courbeC/.style={coulCourbeC,line width=1.1pt,smooth},
}

\newlist{etapes}{enumerate}{1}
\setlist[etapes]{label=\textbf{\arabic*.},leftmargin=*,itemsep=1pt,topsep=2pt}

\newlist{expressions}{enumerate}{1}
\setlist[expressions]{label={\Alph*\ $=$},leftmargin=*,itemsep=3pt,topsep=3pt}

\newcounter{exo}
\renewcommand{\theexo}{\ifnum\value{exo}<10 0\fi\arabic{exo}}
\newtcolorbox{exercice}[1][]{%
  enhanced,breakable,
  colback=white,colframe=coulExo,
  boxrule=0.9pt,arc=2pt,
  left=8pt,right=8pt,top=8pt,bottom=6pt,
  coltitle=white,fonttitle=\bfseries,
  attach boxed title to top left={xshift=8pt,yshift=-\tcboxedtitleheight/2},
  boxed title style={colback=coulExo,arc=2pt,boxrule=0pt},
  phantom={\signet{2}{Exercice \arabic{exo}}},
  title={Exercice \theexo\ifx\relax#1\relax\else{} --- #1\fi}}
\newenvironment{exo}[1][\relax]{\refstepcounter{exo}\begin{exercice}[#1]}{\end{exercice}}

  \newenvironment{correction}{\par}{\par}
  \newcommand{\autableau}{}
  \newcommand{\etapeexemples}{\etape{Exemples corrigés}}
  \newcommand{\etapetableau}[1]{\etape{#1}}
  \newenvironment{exempletableau}[1][Exemple]%
    {\begin{exemplebox}[#1]}{\end{exemplebox}}

\RequirePackage{comment}
\excludecomment{correctionavj}

\newcommand{\coursEntete}{}
\newcommand{\coursNiveau}{}
\newcommand{\coursTitre}{}
\newcommand{\coursSurTitre}{}
\newcommand{\coursSousTitre}{}

\newcommand{\enteteCours}[1]{\renewcommand{\coursEntete}{#1}}
\newcommand{\chapitrecours}[2]{\enteteCours{Chapitre #1 --- #2}}
\newcommand{\niveaucours}[1]{\renewcommand{\coursNiveau}{#1}}
\newcommand{\titrecours}[3][]{%
  \renewcommand{\coursTitre}{#2}%
  \renewcommand{\coursSurTitre}{#1}%
  \renewcommand{\coursSousTitre}{#3}}

\pagestyle{fancy}
\fancyhf{}
\fancyhead[L]{\footnotesize\color{coulPart}\textbf{\coursEntete}}
\fancyhead[R]{\footnotesize\color{coulPart}\textbf{\coursNiveau}}
\fancyfoot[C]{\footnotesize\thepage}
\renewcommand{\headrulewidth}{0.4pt}
\renewcommand{\headrule}{\hbox to\headwidth{\color{coulPart!40}\leaders\hrule height \headrulewidth\hfill}}

\setlength{\parskip}{2pt}

\AtBeginDocument{%
  \begin{center}
    \begin{tcolorbox}[enhanced,width=0.72\linewidth,colback=white,colframe=coulPart,
        boxrule=1.6pt,arc=3pt,halign=center,top=10pt,bottom=10pt]
      \ifx\coursSurTitre\empty
        {\Huge\bfseries\color{coulPart} \coursTitre}\\[4pt]
      \else
        {\Huge\bfseries\color{coulPart} \coursTitre}\\[3pt]
        {\Large\bfseries\color{coulPart} \coursSurTitre}\\[5pt]
      \fi
      {\normalsize \coursSousTitre}%
      
    \end{tcolorbox}
  \end{center}%
}
\makeatother

\chapitrecours{6}{Ensembles de nombres et racines carrées}
\niveaucours{Seconde}
\titrecours{ENSEMBLES DE NOMBRES ET RACINES CARRÉES}{Cours --- Seconde}

\begin{document}

\begin{remarquebox}[Comment utiliser ce cours]
Chaque partie commence par ce qu'il faut \textbf{savoir} (définitions, propriétés), puis
vient ce qu'il faut \textbf{savoir faire} : une méthode, avec des
\textbf{exemples corrigés} faits en classe, puis des exercices
\og\textbf{À vous de jouer}\fg{} à chercher seul.

\end{remarquebox}

\partiecours{Les ensembles de nombres}

\begin{definitionbox}[Définition 1 --- Les ensembles de nombres]
\begin{itemize}
  \item Les \textbf{entiers naturels} servent à compter : $0$ ; $1$ ; $2$ ; $3$\ldots{}
    Leur ensemble est noté $\mathbb{N}$.
  \item Les \textbf{entiers relatifs} sont les entiers positifs ou négatifs : \ldots{} ;
    $-2$ ; $-1$ ; $0$ ; $1$ ; $2$\ldots{} Leur ensemble est noté $\mathbb{Z}$.
  \item Les \textbf{nombres décimaux} s'écrivent avec un nombre \textbf{fini} de chiffres
    après la virgule, c'est-à-dire sous la forme $\dfrac{a}{10^{p}}$ avec $a$ entier
    relatif et $p$ entier naturel : $2{,}07=\dfrac{207}{10^{2}}$. Leur ensemble est
    noté $\mathbb{D}$.
  \item Les \textbf{nombres rationnels} sont les quotients $\dfrac{a}{b}$ de deux entiers
    relatifs, avec $b\ne 0$ : $\dfrac{2}{3}$ ; $-\dfrac{5}{7}$\ldots{} Leur ensemble est
    noté $\mathbb{Q}$.
  \item Les \textbf{nombres réels} sont les abscisses des points de la droite graduée
    (chapitre 2). Leur ensemble est noté $\R$.
\end{itemize}
On écrit $-4\in\mathbb{Z}$ (\og $-4$ appartient à $\mathbb{Z}$\fg) et
$-4\notin\mathbb{N}$ (\og $-4$ n'appartient pas à $\mathbb{N}$\fg).
\end{definitionbox}

\begin{proprietebox}[Propriété 1 --- Les ensembles s'emboîtent]
\begin{minipage}[c]{0.44\linewidth}\raggedright
Un entier naturel est aussi un entier relatif, un entier relatif est aussi un décimal
($-4=-4{,}0$), etc. :
\[
  \mathbb{N}\subset\mathbb{Z}\subset\mathbb{D}\subset\mathbb{Q}\subset\R
\]
Le symbole $\subset$ se lit \og est inclus dans\fg.

\smallskip
$\R$ contient aussi des nombres qui ne sont pas rationnels, comme $\sqrt{2}$ et $\pi$
(partie VII).
\end{minipage}\hfill
\begin{minipage}[c]{0.53\linewidth}\centering
\begin{tikzpicture}[scale=0.82,font=\small]
  \draw[coulDef,thick,rounded corners] (-0.8,-0.8) rectangle (6.2,3.0);
  \draw[coulDef,thick,rounded corners] (-0.6,-0.6) rectangle (5.0,2.5);
  \draw[coulDef,thick,rounded corners] (-0.4,-0.4) rectangle (3.8,2.0);
  \draw[coulDef,thick,rounded corners] (-0.2,-0.2) rectangle (2.6,1.5);
  \draw[coulDef,thick,rounded corners] (0,0) rectangle (1.4,1.0);
  \node at (0.3,0.72) {$\mathbb{N}$};
  \node at (0.1,1.25) {$\mathbb{Z}$};
  \node at (-0.1,1.75) {$\mathbb{D}$};
  \node at (-0.3,2.25) {$\mathbb{Q}$};
  \node at (-0.5,2.75) {$\R$};
  \node at (0.9,0.35) {$7$};
  \node at (2.0,0.35) {$-3$};
  \node at (3.2,0.35) {$0{,}25$};
  \node at (4.4,0.35) {$\frac{1}{3}$};
  \node at (5.6,0.35) {$\sqrt{2}$};
  \node at (5.6,-0.35) {$\pi$};
\end{tikzpicture}
\end{minipage}
\end{proprietebox}

\begin{definitionbox}[Définition 2 --- Nature d'un nombre]
La \textbf{nature} d'un nombre est le \textbf{plus petit} ensemble auquel il appartient.

\smallskip
Exemple : $2$ appartient à $\mathbb{Q}$, car $2=\dfrac{4}{2}$ ; mais $2$ est d'abord un
entier naturel : sa nature est $\mathbb{N}$.
\end{definitionbox}

\begin{remarquebox}[Décimal ou non ?]
Un quotient est décimal quand la division \og tombe juste\fg : $\dfrac{7}{8}=0{,}875$.
Quand elle ne s'arrête jamais, comme $\dfrac{5}{6}=0{,}8333\ldots$, le nombre n'est pas
décimal : il est rationnel, et ses chiffres finissent par se répéter.
\end{remarquebox}

\begin{methodebox}[Méthode 1 --- Trouver la nature d'un nombre]
\etapeexemples
Pour chaque nombre, dire à quels ensembles il appartient, puis donner le plus petit ensemble auquel il appartient (c'est la nature du nombre) :
\begin{questions}[label=\alph*)]
  \item $-\dfrac{18}{6}$
  \item $\dfrac{9}{4}$
  \item $\dfrac{2}{3}$
  \item $\sqrt{49}$
  \item $-0{,}4$
\end{questions}
\begin{correction}
\cas{a)}
\begin{calculs}
-\dfrac{18}{6} &= -3 &\qquad& \rem{$18=3\times 6$ : la division tombe juste}
\end{calculs}
\explique{$-3\notin\mathbb{N}$ et $-3\in\mathbb{Z}$, donc
$-3\in\mathbb{D}$, $-3\in\mathbb{Q}$ et $-3\in\R$.}{un entier négatif n'est pas naturel ;
les ensembles s'emboîtent}
\explique{Sa nature est $\mathbb{Z}$.}{le plus petit ensemble auquel il appartient}
\cas{b)}
\begin{calculs}
\dfrac{9}{4} &= 2{,}25 &\qquad& \rem{la division tombe juste}
\end{calculs}
\explique{$2{,}25$ n'est pas entier : $\dfrac{9}{4}\in\mathbb{D}$, donc aussi à $\mathbb{Q}$
et à $\R$.}{deux chiffres après la virgule : un nombre fini}
\explique{Sa nature est $\mathbb{D}$.}{le plus petit ensemble auquel il appartient}
\cas{c)}
\begin{calculs}
\dfrac{2}{3} &= 0{,}666\ldots &\qquad& \rem{la division ne s'arrête jamais}
\end{calculs}
\explique{$\dfrac{2}{3}\notin\mathbb{D}$, mais $\dfrac{2}{3}\in\mathbb{Q}$ et
$\dfrac{2}{3}\in\R$.}{c'est un quotient de deux entiers}
\explique{Sa nature est $\mathbb{Q}$.}{le plus petit ensemble auquel il appartient}
\cas{d)}
\begin{calculs}
\sqrt{49} &= 7 &\qquad& \rem{$7^{2}=49$ et $7$ est positif}
\end{calculs}
\explique{$7\in\mathbb{N}$, donc $\sqrt{49}$ appartient aux cinq ensembles.}{ne pas se
fier au symbole $\sqrt{\phantom{a}}$ : on calcule d'abord}
\explique{Sa nature est $\mathbb{N}$.}{le plus petit ensemble auquel il appartient}
\cas{e)}
\explique{$-0{,}4$ a un seul chiffre après la virgule : $-0{,}4\in\mathbb{D}$, donc
$-0{,}4\in\mathbb{Q}$ et $-0{,}4\in\R$ ; $-0{,}4\notin\mathbb{Z}$.}{un décimal peut être
négatif}
\explique{Sa nature est $\mathbb{D}$.}{le plus petit ensemble auquel il appartient}
\end{correction}

\etape{À vous de jouer}
Pour chaque nombre, dire à quels ensembles il appartient, puis donner le plus petit ensemble auquel il appartient (c'est la nature du nombre) :
\[
  -\dfrac{35}{7} \qquad \dfrac{13}{5} \qquad \dfrac{4}{9} \qquad \sqrt{81} \qquad -1{,}25
\]
\reponse{$-5$ : $\mathbb{Z}$ ;\; $2{,}6$ : $\mathbb{D}$ ;\; $0{,}444\ldots$ : $\mathbb{Q}$ ;\;
$9$ : $\mathbb{N}$ ;\; $-1{,}25$ : $\mathbb{D}$.}
\end{methodebox}

\afaire{exercices 1 à 4}

\partiecours{Multiples et diviseurs}

\begin{definitionbox}[Définition 3 --- Multiple, diviseur, fraction irréductible]
Soit $a$ et $b$ deux entiers. On dit que $a$ est un \textbf{multiple} de $b$ s'il existe
un entier $k$ tel que $a=k\times b$. On dit aussi que $b$ est un \textbf{diviseur} de $a$,
ou que $a$ est \textbf{divisible} par $b$.

\smallskip
Exemple : $42=6\times 7$, donc $42$ est un multiple de $7$ ; $7$ et $6$ sont des diviseurs
de $42$.

\smallskip
Une fraction est \textbf{irréductible} quand son numérateur et son dénominateur n'ont pas
d'autre diviseur commun positif que $1$.
\end{definitionbox}

\begin{remarquebox}[Rappel du collège : critères de divisibilité]
Un entier est divisible :
\begin{itemize}
  \item par $2$ si son chiffre des unités est $0$ ; $2$ ; $4$ ; $6$ ou $8$ ;
  \item par $5$ s'il se termine par $0$ ou $5$ ;
  \item par $10$ s'il se termine par $0$ ;
  \item par $3$ si la somme de ses chiffres est un multiple de $3$ ;
  \item par $9$ si la somme de ses chiffres est un multiple de $9$.
\end{itemize}
\end{remarquebox}

\begin{methodebox}[Méthode 2 --- Reconnaître un multiple, chercher des diviseurs]
\etapeexemples
\begin{questions}[label=\alph*)]
  \item $126$ est-il un multiple de $7$ ?
  \item $200$ est-il un multiple de $7$ ? Donner le plus grand multiple de $7$ inférieur
    ou égal à $200$.
  \item Donner tous les diviseurs de $45$.
  \item Donner les diviseurs communs à $45$ et $60$, puis rendre la fraction
    $\dfrac{45}{60}$ irréductible.
\end{questions}
\begin{correction}
\cas{a)}
\begin{calculs}
126\div 7 &= 18 &\qquad& \rem{la division tombe juste}
\end{calculs}
\explique{$126=7\times 18$ avec $18$ entier : $126$ est un multiple de $7$.}{on revient à
la définition}
\cas{b)}
\begin{calculs}
200\div 7 &\approx 28{,}57 &\qquad& \rem{le quotient n'est pas entier}
\end{calculs}
\explique{$200$ n'est pas un multiple de $7$.}{}
\begin{calculs}
7\times 28 &= 196 &\qquad& \rem{on garde la partie entière du quotient, $28$}\\
7\times 29 &= 203 &\qquad& \rem{le multiple suivant dépasse $200$}
\end{calculs}
\explique{Le plus grand multiple de $7$ inférieur ou égal à $200$ est $196$.}{}
\cas{c)}
\begin{calculs}
45 &= 1\times 45 &\qquad& \rem{on cherche les diviseurs par paires}\\
45 &= 3\times 15 &\qquad& \rem{$45$ est impair : $2$ ne le divise pas}\\
45 &= 5\times 9 &\qquad& \rem{$4$ ne divise pas $45$}
\end{calculs}
\explique{$6$ ne divise pas $45$, et $7\times 7=49$ dépasse $45$ : on s'arrête.}{au-delà,
on retrouverait les mêmes paires}
\explique{Les diviseurs de $45$ sont $\ens{1\,;3\,;5\,;9\,;15\,;45}$.}{on les range dans
l'ordre}
\cas{d)}
\begin{calculs}
60 &= 1\times 60 &\qquad& \rem{mêmes paires pour $60$}\\
60 &= 2\times 30 &&\\
60 &= 3\times 20 &&\\
60 &= 4\times 15 &&\\
60 &= 5\times 12 &&\\
60 &= 6\times 10 &\qquad& \rem{$7$ ne divise pas $60$ et $8\times 8=64$ dépasse $60$}
\end{calculs}
\explique{Les diviseurs communs à $45$ et $60$ sont $\ens{1\,;3\,;5\,;15}$.}{on garde ceux
des deux listes}
\begin{calculs}
\dfrac{45}{60} &= \dfrac{3\times\barre{15}}{4\times\barre{15}} &\qquad& \rem{on fait apparaître le plus grand diviseur commun, $15$}\\
\dfrac{45}{60} &= \dfrac{3}{4} &\qquad& \rem{$3$ et $4$ n'ont que $1$ comme diviseur commun}
\end{calculs}
\end{correction}

\etape{À vous de jouer}
\begin{questions}[label=\alph*)]
  \item $154$ est-il un multiple de $7$ ?
  \item $150$ est-il un multiple de $7$ ? Donner le plus grand multiple de $7$ inférieur
    ou égal à $150$.
  \item Donner tous les diviseurs de $28$.
  \item Donner les diviseurs communs à $28$ et $42$, puis rendre la fraction
    $\dfrac{28}{42}$ irréductible.
\end{questions}
\reponse{a) oui, $154=7\times 22$ ;\; b) non ; $147$ ;\;
c) $\ens{1\,;2\,;4\,;7\,;14\,;28}$ ;\; d) $\ens{1\,;2\,;7\,;14}$ et
$\dfrac{28}{42}=\dfrac{2}{3}$.}
\end{methodebox}

\afaire{exercices 5 à 9}

\partiecours{Démontrer avec des multiples}

\begin{proprietebox}[Propriété 2 --- Somme de deux multiples]
La somme de deux multiples de $5$ est un multiple de $5$. C'est vrai pour tout entier à la
place de $5$.

\smallskip
\demo{Soit $m$ et $p$ deux multiples de $5$ : il existe deux entiers $k$ et $k'$ tels que
$m=5k$ et $p=5k'$.}
\begin{calculs}
m+p &= 5k+5k' &\qquad& \rem{on remplace $m$ et $p$}\\
m+p &= 5(k+k') &\qquad& \rem{on factorise par $5$}\\
m+p &= 5v &\qquad& \rem{on pose $v=k+k'$, une nouvelle lettre}
\end{calculs}
\explique{$v$ est un entier (somme de deux entiers).}{on revient à la définition :
$m+p=5\times v$}
\explique{Conclusion : $m+p$ est un multiple de $5$.}{}
\end{proprietebox}

\begin{methodebox}[Méthode 3 --- Démontrer qu'un nombre est un multiple]
\etapeexemples
Dans cet exemple, $n$ désigne un entier.
\begin{questions}[label=\alph*)]
  \item Tester pour $n=1$ et $n=2$, puis démontrer que $12n+18$ est un multiple de $6$.
  \item Tester sur deux exemples, puis démontrer que la somme de trois entiers
    consécutifs est un multiple de $3$.
\end{questions}
\begin{correction}
\cas{a)}
\explique{Test : pour $n=1$ : $12\times 1+18=30$ et $30=6\times 5$ ; pour $n=2$ : $12\times 2+18=42$ et $42=6\times 7$.}{on regarde d'abord des exemples}
\begin{calculs}
12n+18 &= 6\times 2n+6\times 3 &\qquad& \rem{on fait apparaître $6$ dans chaque terme}\\
12n+18 &= 6(2n+3) &\qquad& \rem{on factorise par $6$}\\
12n+18 &= 6v &\qquad& \rem{on pose $v=2n+3$}
\end{calculs}
\explique{$v$ est un entier.}{}
\explique{Conclusion : $12n+18$ est un multiple de $6$.}{on revient à la définition}
\cas{b)}
\explique{Test : $4+5+6=15$ et $15=3\times 5$ ; $10+11+12=33$ et $33=3\times 11$.}{on regarde d'abord des exemples}
\explique{Trois entiers consécutifs s'écrivent $n$, $n+1$ et $n+2$.}{chacun vaut $1$ de
plus que le précédent}
\begin{calculs}
S &= n+(n+1)+(n+2) &\qquad& \rem{on nomme la somme $S$}\\
S &= 3n+3 &\qquad& \rem{on réduit}\\
S &= 3(n+1) &\qquad& \rem{on factorise par $3$}\\
S &= 3v &\qquad& \rem{on pose $v=n+1$}
\end{calculs}
\explique{$v$ est un entier.}{}
\explique{Conclusion : la somme de trois entiers consécutifs est un multiple de $3$.}{}
\end{correction}

\etape{À vous de jouer}
$n$ désigne un entier.
\begin{questions}[label=\alph*)]
  \item Tester pour $n=1$ et $n=2$, puis démontrer que $15n+35$ est un multiple de $5$.
  \item Tester sur deux exemples, puis démontrer que la somme de cinq entiers
    consécutifs est un multiple de $5$.
\end{questions}
\reponse{a) $15n+35=5v$ avec $v=3n+7$ ;\; b) $S=5v$ avec $v=n+2$.}
\end{methodebox}

\begin{remarquebox}[Montrer qu'une affirmation est fausse]
\og La somme de quatre entiers consécutifs est un multiple de $4$\fg{} est fausse :
$1+2+3+4=10$, qui n'est pas un multiple de $4$. Un seul \textbf{contre-exemple} suffit pour
montrer qu'une affirmation est fausse ; des exemples ne suffisent jamais pour montrer
qu'elle est vraie.
\end{remarquebox}

\afaire{exercices 10 à 12}

\partiecours{Nombres pairs et impairs}

\begin{definitionbox}[Définition 4 --- Entier pair, entier impair]
Un entier \textbf{pair} est un multiple de $2$ : il s'écrit $2k$, avec $k$ entier.

\smallskip
Un entier \textbf{impair} est un entier qui n'est pas pair : il s'écrit $2k+1$, avec $k$
entier.

\smallskip
Exemples : $34=2\times 17$ est pair ; $51=2\times 25+1$ est impair ; $0=2\times 0$ est pair.
\end{definitionbox}

\begin{methodebox}[Méthode 4 --- Écrire un nombre sous la forme $2k$ ou $2k+1$]
\etapeexemples
Dans cet exemple, $n$ désigne un entier.
\begin{questions}[label=\alph*)]
  \item Écrire $86$ et $123$ sous la forme $2k$ ou $2k+1$.
  \item Tester pour $n=1$ et $n=2$, puis démontrer que $8n+6$ est pair.
  \item Tester pour $n=1$ et $n=2$, puis démontrer que $6n+11$ est impair.
\end{questions}
\begin{correction}
\cas{a)}
\begin{calculs}
86 &= 2\times 43 &\qquad& \rem{$86$ est pair : $k=43$}\\
123 &= 122+1 &\qquad& \rem{on fait apparaître le pair qui précède}\\
123 &= 2\times 61+1 &\qquad& \rem{$123$ est impair : $k=61$}
\end{calculs}
\cas{b)}
\explique{Test : pour $n=1$ : $14=2\times 7$ ; pour $n=2$ : $22=2\times 11$.}{on regarde d'abord des exemples}
\begin{calculs}
8n+6 &= 2(4n+3) &\qquad& \rem{on factorise par $2$}\\
8n+6 &= 2v &\qquad& \rem{on pose $v=4n+3$}
\end{calculs}
\explique{$v$ est un entier.}{forme $2k$, avec la lettre $v$}
\explique{Conclusion : $8n+6$ est pair.}{}
\cas{c)}
\explique{Test : pour $n=1$ : $17=2\times 8+1$ ; pour $n=2$ : $23=2\times 11+1$.}{on regarde d'abord des exemples}
\begin{calculs}
6n+11 &= 6n+10+1 &\qquad& \rem{on met de côté $1$ pour que le reste soit pair}\\
6n+11 &= 2(3n+5)+1 &\qquad& \rem{on factorise $6n+10$ par $2$}\\
6n+11 &= 2v+1 &\qquad& \rem{on pose $v=3n+5$}
\end{calculs}
\explique{$v$ est un entier.}{forme $2k+1$, avec la lettre $v$}
\explique{Conclusion : $6n+11$ est impair.}{}
\end{correction}

\etape{À vous de jouer}
$n$ désigne un entier.
\begin{questions}[label=\alph*)]
  \item Écrire $74$ et $145$ sous la forme $2k$ ou $2k+1$.
  \item Tester pour $n=1$ et $n=2$, puis démontrer que $10n+4$ est pair.
  \item Tester pour $n=1$ et $n=2$, puis démontrer que $4n+9$ est impair.
\end{questions}
\reponse{a) $74=2\times 37$ ;\; $145=2\times 72+1$ ;\; b) $10n+4=2v$ avec $v=5n+2$ ;\;
c) $4n+9=2v+1$ avec $v=2n+4$.}
\end{methodebox}

\begin{proprietebox}[Propriété 3 --- Parité d'une somme]
\begin{itemize}
  \item pair $+$ pair donne un nombre pair ;
  \item pair $+$ impair donne un nombre impair ;
  \item impair $+$ impair donne un nombre pair.
\end{itemize}
Les deux dernières se démontrent dans la méthode 5.
\end{proprietebox}

\begin{methodebox}[Méthode 5 --- Démontrer une règle sur deux entiers]
\etapeexemples
\begin{questions}[label=\alph*)]
  \item Tester sur deux exemples, puis démontrer que la somme de deux entiers impairs
    est paire.
  \item Tester sur deux exemples, puis démontrer que la somme d'un entier pair et d'un
    entier impair est impaire.
\end{questions}
\begin{correction}
\cas{a)}
\explique{Test : $3+7=10$ et $5+9=14$ sont pairs.}{on regarde d'abord des exemples}
\explique{Soit $a$ et $b$ deux entiers impairs : $a=2k+1$ et $b=2k'+1$, avec $k$ et $k'$
entiers.}{deux lettres différentes : $a$ et $b$ ne sont pas forcément égaux}
\begin{calculs}
a+b &= 2k+1+2k'+1 &\qquad& \rem{on remplace $a$ et $b$}\\
a+b &= 2k+2k'+2 &\qquad& \rem{on réduit}\\
a+b &= 2(k+k'+1) &\qquad& \rem{on factorise par $2$}\\
a+b &= 2v &\qquad& \rem{on pose $v=k+k'+1$}
\end{calculs}
\explique{$v$ est un entier.}{forme $2k$}
\explique{Conclusion : la somme de deux entiers impairs est paire.}{}
\cas{b)}
\explique{Test : $4+7=11$ et $10+3=13$ sont impairs.}{on regarde d'abord des exemples}
\explique{Soit $a$ un entier pair et $b$ un entier impair : $a=2k$ et $b=2k'+1$, avec $k$
et $k'$ entiers.}{}
\begin{calculs}
a+b &= 2k+2k'+1 &\qquad& \rem{on remplace $a$ et $b$}\\
a+b &= 2(k+k')+1 &\qquad& \rem{on factorise les deux premiers termes par $2$}\\
a+b &= 2v+1 &\qquad& \rem{on pose $v=k+k'$}
\end{calculs}
\explique{$v$ est un entier.}{forme $2k+1$}
\explique{Conclusion : la somme d'un entier pair et d'un entier impair est impaire.}{}
\end{correction}

\etape{À vous de jouer}
\begin{questions}[label=\alph*)]
  \item Tester sur deux exemples, puis démontrer que la différence de deux entiers
    impairs est paire.
  \item Tester sur deux exemples, puis démontrer que la somme de deux entiers pairs est
    paire.
\end{questions}
\reponse{a) $a-b=2v$ avec $v=k-k'$ ;\; b) $a+b=2v$ avec $v=k+k'$.}
\end{methodebox}

\afaire{exercices 13 à 16}

\partiecours{Démonstrations : nombres pairs et impairs}

\begin{proprietebox}[Propriété 4 --- Carré d'un entier impair]
Le carré d'un entier impair est impair.

\smallskip
\demo{Soit $a$ un entier impair : $a=2k+1$, avec $k$ entier.}
\begin{calculs}
a^{2} &= (2k+1)^{2} &\qquad& \rem{on remplace $a$}\\
a^{2} &= 4k^{2}+4k+1 &\qquad& \rem{identité remarquable $(a+b)^{2}$ (chapitre 4)}\\
a^{2} &= 2(2k^{2}+2k)+1 &\qquad& \rem{on factorise $4k^{2}+4k$ par $2$}\\
a^{2} &= 2v+1 &\qquad& \rem{on pose $v=2k^{2}+2k$}
\end{calculs}
\explique{$v$ est un entier.}{forme $2k+1$}
\explique{Conclusion : le carré d'un entier impair est impair.}{}
\end{proprietebox}

\begin{remarquebox}[Conséquence, utile à la partie VII]
Si $a^{2}$ est pair, alors $a$ est pair : si $a$ était impair, $a^{2}$ serait impair
(propriété 4).
\end{remarquebox}

\begin{methodebox}[Méthode 6 --- Démontrer en séparant les cas pair et impair]
\etapeexemples
On ne sait pas si l'entier $n$ est pair ou impair : on traite les deux cas
(\emph{disjonction des cas}). Tester pour $n=2$ et $n=3$, puis démontrer que, pour tout
entier $n$, $n^{2}+n$ est pair.
\begin{correction}
\cas{Résolution}\par\nopagebreak
\explique{Test : pour $n=2$ : $4+2=6$ ; pour $n=3$ : $9+3=12$. Les deux sont pairs.}{un entier pair, un entier impair : les deux cas}
\explique{\textbf{1\textsuperscript{er} cas} : $n$ est pair, $n=2k$ avec $k$ entier.}{}
\begin{calculs}
n^{2}+n &= (2k)^{2}+2k &\qquad& \rem{on remplace $n$}\\
n^{2}+n &= 4k^{2}+2k &&\\
n^{2}+n &= 2(2k^{2}+k) &\qquad& \rem{on factorise par $2$}\\
n^{2}+n &= 2v &\qquad& \rem{on pose $v=2k^{2}+k$, entier : pair}
\end{calculs}
\explique{\textbf{2\textsuperscript{e} cas} : $n$ est impair, $n=2k+1$ avec $k$ entier.}{}
\begin{calculs}
n^{2}+n &= (2k+1)^{2}+2k+1 &\qquad& \rem{on remplace $n$}\\
n^{2}+n &= 4k^{2}+4k+1+2k+1 &\qquad& \rem{identité remarquable}\\
n^{2}+n &= 4k^{2}+6k+2 &\qquad& \rem{on réduit}\\
n^{2}+n &= 2(2k^{2}+3k+1) &\qquad& \rem{on factorise par $2$}\\
n^{2}+n &= 2w &\qquad& \rem{on pose $w=2k^{2}+3k+1$, entier : pair}
\end{calculs}
\explique{Conclusion : dans les deux cas, $n^{2}+n$ est pair : il est pair pour tout entier $n$.}{on conclut sur tous les entiers}
\end{correction}

\etape{À vous de jouer}
Tester pour $n=2$ et $n=3$, puis démontrer que, pour tout entier $n$ :
\begin{questions}[label=\alph*)]
  \item $n^{2}+5n$ est pair ;
  \item $3n^{2}+n$ est pair.
\end{questions}
\reponse{a) $2v$ avec $v=2k^{2}+5k$, puis $2w$ avec $w=2k^{2}+7k+3$ ;\;
b) $2v$ avec $v=6k^{2}+k$, puis $2w$ avec $w=6k^{2}+7k+2$.}
\end{methodebox}

\afaire{exercices 17 à 20}

\partiecours{Démonstration par l'absurde : $\frac{1}{3}$ n'est pas décimal}

\begin{definitionbox}[Définition 5 --- Raisonnement par l'absurde]
Pour démontrer qu'une affirmation est vraie, on suppose qu'elle est \textbf{fausse}, et on
en déduit quelque chose d'\textbf{impossible}. La supposition était donc fausse :
l'affirmation est vraie.

\smallskip
Exemple : la vitre de la salle a été cassée à 7 h 30. Supposons que Mathéo l'ait cassée.
Or, à 7 h 30, la caméra du bus le filme en route vers le lycée : c'est impossible.
Conclusion : Mathéo n'a pas cassé la vitre.
\end{definitionbox}

\begin{proprietebox}[Propriété 5 --- $\frac{1}{3}$ n'est pas décimal]
Le nombre rationnel $\dfrac{1}{3}$ n'est pas décimal.

\smallskip
\demo{Par l'absurde, supposons que $\dfrac{1}{3}$ soit décimal : $\dfrac{1}{3}=\dfrac{a}{10^{p}}$
avec $a$ entier et $p$ entier naturel.}
\begin{calculs}
\dfrac{1}{3} &= \dfrac{a}{10^{p}} &\qquad& \rem{forme d'un nombre décimal}\\
10^{p} &= 3a &\qquad& \rem{produit en croix (chapitre 1)}
\end{calculs}
\explique{Donc $10^{p}$ est un multiple de $3$.}{$a$ est entier}
\explique{Or $10^{p}$ s'écrit avec un $1$ suivi de zéros.}{$10^{3}=1000$}
\explique{La somme de ses chiffres vaut $1$.}{}
\explique{$1$ n'est pas un multiple de $3$.}{}
\explique{Donc $10^{p}$ n'est pas un multiple de $3$.}{critère de divisibilité par $3$}
\explique{C'est impossible.}{la supposition était fausse}
\explique{Conclusion : $\dfrac{1}{3}$ n'est pas décimal.}{}
\end{proprietebox}

\begin{methodebox}[Méthode 7 --- Démontrer par l'absurde]
\etapeexemples
\begin{questions}[label=\alph*)]
  \item Démontrer que $\dfrac{2}{3}$ n'est pas décimal.
  \item Démontrer qu'il n'existe pas de plus grand entier pair.
\end{questions}
\begin{correction}
\cas{a)}
\explique{Supposons que $\dfrac{2}{3}$ soit décimal : $\dfrac{2}{3}=\dfrac{a}{10^{p}}$ avec
$a$ entier et $p$ entier naturel.}{on suppose le contraire}
\begin{calculs}
\dfrac{2}{3} &= \dfrac{a}{10^{p}} &&\\
2\times 10^{p} &= 3a &\qquad& \rem{produit en croix}
\end{calculs}
\explique{Donc $2\times 10^{p}$ est un multiple de $3$.}{}
\explique{Or $2\times 10^{p}$ s'écrit avec un $2$ suivi de zéros.}{$2\times 10^{3}=2000$}
\explique{La somme de ses chiffres vaut $2$.}{}
\explique{$2$ n'est pas un multiple de $3$.}{}
\explique{Donc $2\times 10^{p}$ n'est pas un multiple de $3$.}{critère de divisibilité par
$3$}
\explique{C'est impossible.}{la supposition était fausse}
\explique{Conclusion : $\dfrac{2}{3}$ n'est pas décimal.}{}
\cas{b)}
\explique{Supposons qu'il existe un plus grand entier pair $N$ : $N=2k$, avec $k$ entier.}{on
suppose le contraire}
\begin{calculs}
N+2 &= 2k+2 &\qquad& \rem{on regarde l'entier pair suivant}\\
N+2 &= 2(k+1) &\qquad& \rem{on factorise par $2$}\\
N+2 &= 2v &\qquad& \rem{on pose $v=k+1$, entier : $N+2$ est pair}
\end{calculs}
\explique{$N+2$ est pair et plus grand que $N$.}{}
\explique{C'est impossible, puisque $N$ est le plus grand.}{on a trouvé une
contradiction}
\explique{Conclusion : il n'existe pas de plus grand entier pair.}{}
\end{correction}

\etape{À vous de jouer}
\begin{questions}[label=\alph*)]
  \item Démontrer que $\dfrac{1}{9}$ n'est pas décimal.
  \item Démontrer qu'il n'existe pas de plus grand multiple de $5$.
\end{questions}
\reponse{a) $10^{p}=9a$, mais la somme des chiffres de $10^{p}$ vaut $1$ ;\;
b) $N+5=5v$ avec $v=k+1$ est un multiple de $5$ plus grand que $N$.}
\end{methodebox}

\afaire{exercices 21 à 23}

\partiecours{Nombres irrationnels : $\sqrt{2}$ n'est pas rationnel}

\begin{definitionbox}[Définition 6 --- Nombre irrationnel]
\begin{minipage}[c]{0.58\linewidth}\raggedright
Un nombre réel qui n'est pas rationnel est dit \textbf{irrationnel}.

\smallskip
La géométrie en fournit : la diagonale d'un carré de côté $1$ mesure $\sqrt{2}$
(théorème de Pythagore : $1^{2}+1^{2}=2$) ; le périmètre d'un cercle de diamètre $1$
mesure $\pi$. Ces deux nombres sont irrationnels ; pour $\pi$, on l'admet.

\smallskip
Un nombre irrationnel a une infinité de chiffres après la virgule, sans répétition :
la calculatrice n'en affiche qu'une \textbf{valeur approchée}.
\end{minipage}\hfill
\begin{minipage}[c]{0.38\linewidth}\centering
\begin{tikzpicture}[x=1.8cm,y=1.8cm,font=\small]
  \draw[->] (-0.3,0)--(1.8,0);
  \draw[coulDef,thick] (0,0) rectangle (1,1);
  \draw[coulDef,thick] (0,0)--node[above left=-1pt,font=\footnotesize]{$\sqrt{2}$} (1,1);
  \draw[dashed] (1,1) arc[start angle=45,end angle=0,radius=1.41421];
  \path (0,0) node[croix]{} node[below=2pt]{$0$};
  \path (1,0) node[croix]{} node[below=2pt]{$1$};
  \path (1.41421,0) node[croix]{} node[below=2pt]{$\sqrt{2}$};
  \node[left] at (0,0.5) {$1$};
\end{tikzpicture}
\end{minipage}
\end{definitionbox}

Pour la démonstration, on se sert de deux résultats déjà vus : toute fraction peut être
rendue \textbf{irréductible} (partie II), et si $a^{2}$ est pair, alors $a$ est pair
(partie V).

\begin{proprietebox}[Propriété 6 --- $\sqrt{2}$ est irrationnel]
Le nombre $\sqrt{2}$ n'est pas rationnel.

\smallskip
\demo{Par l'absurde, supposons que $\sqrt{2}$ soit rationnel : $\sqrt{2}=\dfrac{a}{b}$, où
$\dfrac{a}{b}$ est une fraction irréductible d'entiers positifs.}
\begin{calculs}
\sqrt{2} &= \dfrac{a}{b} &&\\
2 &= \dfrac{a^{2}}{b^{2}} &\qquad& \rem{on élève au carré : $\big(\sqrt{2}\big)^{2}=2$}\\
a^{2} &= 2b^{2} &\qquad& \rem{produit en croix}
\end{calculs}
\explique{$a^{2}$ est pair, donc $a$ est pair : $a=2k$, avec $k$ entier.}{conséquence de la
propriété 4}
\begin{calculs}
(2k)^{2} &= 2b^{2} &\qquad& \rem{on remplace $a$}\\
4k^{2} &= 2b^{2} &&\\
b^{2} &= 2k^{2} &\qquad& \rem{on divise par $2$}
\end{calculs}
\explique{$b^{2}$ est pair, donc $b$ est pair.}{même raisonnement}
\explique{$a$ et $b$ sont pairs : $2$ est un diviseur commun, alors que $\dfrac{a}{b}$ est
irréductible. C'est impossible.}{la supposition était fausse}
\explique{Conclusion : $\sqrt{2}$ est irrationnel.}{}
\end{proprietebox}

Pour $\sqrt{3}$, on remplace \og pair\fg{} par \og multiple de $3$\fg. Divisé par $3$, un entier
a pour reste $0$ ; $1$ ou $2$ : il s'écrit $3k$ ; $3k+1$ ou $3k+2$, avec $k$ entier
($17=3\times 5+2$).

\begin{methodebox}[Méthode 8 --- Démontrer qu'une racine carrée est irrationnelle]
\etapeexemples
\begin{questions}[label=\alph*)]
  \item Démontrer que, si $a^{2}$ est un multiple de $3$, alors $a$ est un multiple de $3$.
  \item En déduire que $\sqrt{3}$ est irrationnel.
\end{questions}
\begin{correction}
\cas{a)}
\explique{Supposons que $a$ ne soit pas un multiple de $3$ : $a=3k+1$ ou $a=3k+2$, avec $k$
entier.}{on suppose le contraire, puis on sépare les deux cas}
\begin{calculs}
a^{2} &= (3k+1)^{2} &\qquad& \rem{1\textsuperscript{er} cas}\\
a^{2} &= 9k^{2}+6k+1 &\qquad& \rem{identité remarquable}\\
a^{2} &= 3(3k^{2}+2k)+1 &\qquad& \rem{on factorise $9k^{2}+6k$ par $3$}\\
a^{2} &= 3v+1 &\qquad& \rem{on pose $v=3k^{2}+2k$}\\[4pt]
a^{2} &= (3k+2)^{2} &\qquad& \rem{2\textsuperscript{e} cas}\\
a^{2} &= 9k^{2}+12k+4 &&\\
a^{2} &= 3(3k^{2}+4k+1)+1 &\qquad& \rem{$4=3+1$}\\
a^{2} &= 3w+1 &\qquad& \rem{on pose $w=3k^{2}+4k+1$}
\end{calculs}
\explique{Dans les deux cas, $a^{2}$ a pour reste $1$ dans la division par $3$ : $a^{2}$
n'est pas un multiple de $3$.}{$3v<3v+1<3v+3$ : entre deux multiples de $3$ consécutifs}
\explique{C'est impossible, puisque $a^{2}$ est un multiple de $3$.}{la supposition était
fausse}
\explique{Conclusion : si $a^{2}$ est un multiple de $3$, alors $a$ est un multiple de $3$.}{}
\cas{b)}
\explique{Par l'absurde, supposons $\sqrt{3}=\dfrac{a}{b}$, où $\dfrac{a}{b}$ est une
fraction irréductible d'entiers positifs.}{même démarche que pour $\sqrt{2}$}
\begin{calculs}
\sqrt{3} &= \dfrac{a}{b} &&\\
3 &= \dfrac{a^{2}}{b^{2}} &\qquad& \rem{on élève au carré}\\
a^{2} &= 3b^{2} &\qquad& \rem{produit en croix}
\end{calculs}
\explique{$a^{2}$ est un multiple de $3$, donc $a$ aussi : $a=3k$, avec $k$ entier.}{question
a)}
\begin{calculs}
(3k)^{2} &= 3b^{2} &\qquad& \rem{on remplace $a$}\\
9k^{2} &= 3b^{2} &&\\
b^{2} &= 3k^{2} &\qquad& \rem{on divise par $3$}
\end{calculs}
\explique{$b^{2}$ est un multiple de $3$, donc $b$ aussi.}{question a)}
\explique{$3$ est un diviseur commun de $a$ et $b$, alors que $\dfrac{a}{b}$ est
irréductible.}{}
\explique{C'est impossible.}{la supposition était fausse}
\explique{Conclusion : $\sqrt{3}$ est irrationnel.}{}
\end{correction}

\etape{À vous de jouer}
Un entier s'écrit $5k$ ; $5k+1$ ; $5k+2$ ; $5k+3$ ou $5k+4$, avec $k$ entier.
\begin{questions}[label=\alph*)]
  \item Démontrer que, si $a^{2}$ est un multiple de $5$, alors $a$ est un multiple de $5$.
  \item En déduire que $\sqrt{5}$ est irrationnel.
\end{questions}
\reponse{a) les carrés valent $5v+1$ ; $5v+4$ ; $5v+4$ ; $5v+1$ : jamais un multiple de $5$ ;\;
b) $a^{2}=5b^{2}$, puis $b^{2}=5k^{2}$ : $5$ divise $a$ et $b$.}
\end{methodebox}

\afaire{exercices 24 à 27}

\partiecours{Calculer avec les racines carrées}

\begin{definitionbox}[Définition 7 --- Racine carrée]
Soit $a$ un nombre positif. La \textbf{racine carrée} de $a$, notée $\sqrt{a}$, est le nombre
\textbf{positif} dont le carré vaut $a$ : $\big(\sqrt{a}\big)^{2}=a$.

\smallskip
Exemples : $\sqrt{25}=5$, car $5^{2}=25$ et $5$ est positif ; $\sqrt{0}=0$ ; $\sqrt{1}=1$.
$\sqrt{-4}$ n'existe pas : un carré n'est jamais négatif.

\smallskip
Pour $a$ positif, $\sqrt{a^{2}}=a$ : $\sqrt{7^{2}}=7$.
\end{definitionbox}

\begin{proprietebox}[Propriété 7 --- Racine d'un produit, d'un quotient]
Pour $a$ et $b$ positifs : \quad $\sqrt{a\times b}=\sqrt{a}\times\sqrt{b}$ \quad et, si
$b\ne 0$, \quad $\sqrt{\dfrac{a}{b}}=\dfrac{\sqrt{a}}{\sqrt{b}}$.
\end{proprietebox}

\begin{remarquebox}[L'erreur à éviter]
La racine d'une \textbf{somme} n'est pas la somme des racines :
$\sqrt{36+64}=\sqrt{100}=10$, alors que $\sqrt{36}+\sqrt{64}=6+8=14$. Dans
$\sqrt{a+b}$, on calcule d'abord la somme.
\end{remarquebox}

\begin{methodebox}[Méthode 9 --- Calculer avec des racines carrées]
\etapeexemples
Calculer :
\begin{questions}[label=\alph*)]
  \item $A=\sqrt{2}\times\sqrt{32}$
  \item $B=\dfrac{\sqrt{75}}{\sqrt{3}}$
  \item $C=\big(4\sqrt{3}\big)^{2}$
  \item $D=\sqrt{\dfrac{9}{49}}$
  \item $E=\sqrt{144+25}$
\end{questions}
\begin{correction}
\cas{a)}
\begin{calculs}
A &= \sqrt{2}\times\sqrt{32} &&\\
A &= \sqrt{2\times 32} &\qquad& \rem{racine d'un produit}\\
A &= \sqrt{64} &&\\
A &= 8 &\qquad& \rem{$8^{2}=64$}
\end{calculs}
\cas{b)}
\begin{calculs}
B &= \dfrac{\sqrt{75}}{\sqrt{3}} &&\\
B &= \sqrt{\dfrac{75}{3}} &\qquad& \rem{racine d'un quotient}\\
B &= \sqrt{25} &&\\
B &= 5 &&
\end{calculs}
\cas{c)}
\begin{calculs}
C &= \big(4\sqrt{3}\big)^{2} &&\\
C &= 4^{2}\times\big(\sqrt{3}\big)^{2} &\qquad& \rem{puissance d'un produit : chaque facteur au carré}\\
C &= 16\times 3 &\qquad& \rem{$\big(\sqrt{3}\big)^{2}=3$}\\
C &= 48 &&
\end{calculs}
\cas{d)}
\begin{calculs}
D &= \sqrt{\dfrac{9}{49}} &&\\
D &= \dfrac{\sqrt{9}}{\sqrt{49}} &\qquad& \rem{racine d'un quotient}\\
D &= \dfrac{3}{7} &&
\end{calculs}
\cas{e)}
\begin{calculs}
E &= \sqrt{144+25} &&\\
E &= \sqrt{169} &\qquad& \rem{une somme sous la racine : on la calcule d'abord}\\
E &= 13 &\qquad& \rem{$13^{2}=169$ ; et non $12+5=17$}
\end{calculs}
\end{correction}

\etape{À vous de jouer}
Calculer :
\[
  F=\sqrt{3}\times\sqrt{27} \qquad G=\dfrac{\sqrt{98}}{\sqrt{2}} \qquad
  H=\big(2\sqrt{5}\big)^{2} \qquad I=\sqrt{\dfrac{25}{64}} \qquad J=\sqrt{64+225}
\]
\reponse{$F=9$ ;\; $G=7$ ;\; $H=20$ ;\; $I=\dfrac{5}{8}$ ;\; $J=17$.}
\end{methodebox}

\afaire{exercices 28 à 31}

\partiecours{Transformer et réduire des expressions}

\begin{remarquebox}[Les carrés à connaître]
$4$ ; $9$ ; $16$ ; $25$ ; $36$ ; $49$ ; $64$ ; $81$ ; $100$ ; $121$ ; $144$ sont les carrés
de $2$ à $12$ : on les cherche dans le nombre sous la racine.
\end{remarquebox}

\begin{methodebox}[Méthode 10 --- Écrire une racine sous la forme $a\sqrt{b}$]
\etapeexemples
Écrire sous la forme $a\sqrt{b}$, avec $b$ entier le plus petit possible :
\begin{questions}[label=\alph*)]
  \item $\sqrt{18}$
  \item $\sqrt{48}$
  \item $\sqrt{200}$
  \item $3\sqrt{20}$
\end{questions}
\begin{correction}
\cas{a)}
\begin{calculs}
\sqrt{18} &= \sqrt{9\times 2} &\qquad& \rem{$9$ est le plus grand carré qui divise $18$}\\
\sqrt{18} &= \sqrt{9}\times\sqrt{2} &\qquad& \rem{racine d'un produit}\\
\sqrt{18} &= 3\sqrt{2} &&
\end{calculs}
\cas{b)}
\begin{calculs}
\sqrt{48} &= \sqrt{16\times 3} &\qquad& \rem{$16$ divise $48$ ; $36$ non}\\
\sqrt{48} &= \sqrt{16}\times\sqrt{3} &&\\
\sqrt{48} &= 4\sqrt{3} &&
\end{calculs}
\cas{c)}
\begin{calculs}
\sqrt{200} &= \sqrt{100\times 2} &&\\
\sqrt{200} &= \sqrt{100}\times\sqrt{2} &&\\
\sqrt{200} &= 10\sqrt{2} &&
\end{calculs}
\cas{d)}
\begin{calculs}
3\sqrt{20} &= 3\times\sqrt{4\times 5} &\qquad& \rem{on transforme la racine, le $3$ attend}\\
3\sqrt{20} &= 3\times 2\sqrt{5} &\qquad& \rem{$\sqrt{4}=2$}\\
3\sqrt{20} &= 6\sqrt{5} &&
\end{calculs}
\end{correction}

\etape{À vous de jouer}
Écrire sous la forme $a\sqrt{b}$, avec $b$ entier le plus petit possible :
\[
  \sqrt{27} \qquad \sqrt{80} \qquad \sqrt{300} \qquad 2\sqrt{45}
\]
\reponse{$3\sqrt{3}$ ;\; $4\sqrt{5}$ ;\; $10\sqrt{3}$ ;\; $6\sqrt{5}$.}
\end{methodebox}

\begin{remarquebox}[Carré trop petit]
Avec un carré trop petit, on n'a pas fini : $\sqrt{48}=\sqrt{4\times 12}=2\sqrt{12}$, et
$\sqrt{12}$ contient encore le carré $4$. On recommence : $2\sqrt{12}=2\times 2\sqrt{3}$,
soit $4\sqrt{3}$.
\end{remarquebox}

On réduit $5\sqrt{3}+2\sqrt{3}$ comme $5x+2x=7x$ : on regroupe les termes qui ont
\textbf{la même racine}. Mais $\sqrt{2}+\sqrt{3}$ ne se réduit pas, comme $x+y$.

\begin{methodebox}[Méthode 11 --- Réduire une somme de racines carrées]
\etapeexemples
Réduire :
\begin{questions}[label=\alph*)]
  \item $A=4\sqrt{7}-9\sqrt{7}+\sqrt{7}$
  \item $B=\sqrt{12}+5\sqrt{3}$
  \item $C=\sqrt{50}-3\sqrt{8}+\sqrt{18}$
  \item $D=\sqrt{20}+\sqrt{45}-\sqrt{3}$
\end{questions}
\begin{correction}
\cas{a)}
\begin{calculs}
A &= 4\sqrt{7}-9\sqrt{7}+\sqrt{7} &\qquad& \rem{$\sqrt{7}$ compte pour $1\sqrt{7}$}\\
A &= (4-9+1)\sqrt{7} &\qquad& \rem{même racine : on factorise}\\
A &= -4\sqrt{7} &&
\end{calculs}
\cas{b)}
\begin{calculs}
B &= \sqrt{12}+5\sqrt{3} &&\\
B &= 2\sqrt{3}+5\sqrt{3} &\qquad& \rem{$\sqrt{12}=\sqrt{4\times 3}$ (méthode 10)}\\
B &= 7\sqrt{3} &&
\end{calculs}
\cas{c)}
\begin{calculs}
C &= \sqrt{50}-3\sqrt{8}+\sqrt{18} &&\\
C &= 5\sqrt{2}-3\times 2\sqrt{2}+3\sqrt{2} &\qquad& \rem{$50=25\times 2$ ; $8=4\times 2$ ; $18=9\times 2$}\\
C &= 5\sqrt{2}-6\sqrt{2}+3\sqrt{2} &&\\
C &= 2\sqrt{2} &\qquad& \rem{$5-6+3=2$}
\end{calculs}
\cas{d)}
\begin{calculs}
D &= \sqrt{20}+\sqrt{45}-\sqrt{3} &&\\
D &= 2\sqrt{5}+3\sqrt{5}-\sqrt{3} &\qquad& \rem{$20=4\times 5$ ; $45=9\times 5$ ; $3$ ne contient pas de carré}\\
D &= 5\sqrt{5}-\sqrt{3} &\qquad& \rem{$\sqrt{5}$ et $\sqrt{3}$ ne se regroupent pas}
\end{calculs}
\end{correction}

\etape{À vous de jouer}
Réduire :
\[
  E=2\sqrt{5}-7\sqrt{5}+3\sqrt{5} \qquad F=\sqrt{28}+4\sqrt{7} \qquad
  G=\sqrt{75}-2\sqrt{12}+\sqrt{27} \qquad H=\sqrt{8}+\sqrt{32}-\sqrt{5}
\]
\reponse{$E=-2\sqrt{5}$ ;\; $F=6\sqrt{7}$ ;\; $G=4\sqrt{3}$ ;\; $H=6\sqrt{2}-\sqrt{5}$.}
\end{methodebox}

\afaire{exercices 32 à 35}

\partiecours{Simplifier des expressions avec des racines}

Une racine carrée se manipule comme une lettre : on développe avec la distributivité et
les identités remarquables du chapitre 4, puis on remplace $\big(\sqrt{a}\big)^{2}$ par $a$.
\[
  (a+b)^{2}=a^{2}+2ab+b^{2} \qquad (a-b)^{2}=a^{2}-2ab+b^{2} \qquad (a+b)(a-b)=a^{2}-b^{2}
\]

\begin{methodebox}[Méthode 12 --- Développer une expression avec des racines carrées]
\etapeexemples
Développer et réduire :
\begin{questions}[label=\alph*)]
  \item $A=\sqrt{5}\big(3+\sqrt{5}\big)$
  \item $B=\big(2+\sqrt{3}\big)\big(4-\sqrt{3}\big)$
  \item $C=\big(\sqrt{6}+1\big)^{2}$
  \item $D=\big(5-\sqrt{2}\big)\big(5+\sqrt{2}\big)$
\end{questions}
\begin{correction}
\cas{a)}
\begin{calculs}
A &= \sqrt{5}\big(3+\sqrt{5}\big) &&\\
A &= 3\sqrt{5}+\sqrt{5}\times\sqrt{5} &\qquad& \rem{distributivité simple}\\
A &= 3\sqrt{5}+5 &\qquad& \rem{$\sqrt{5}\times\sqrt{5}=5$}
\end{calculs}
\cas{b)}
\begin{calculs}
B &= \big(2+\sqrt{3}\big)\big(4-\sqrt{3}\big) &&\\
B &= 8-2\sqrt{3}+4\sqrt{3}-\sqrt{3}\times\sqrt{3} &\qquad& \rem{double distributivité}\\
B &= 8-2\sqrt{3}+4\sqrt{3}-3 &&\\
B &= 5+2\sqrt{3} &\qquad& \rem{on regroupe les entiers, puis les $\sqrt{3}$}
\end{calculs}
\cas{c)}
\begin{calculs}
C &= \big(\sqrt{6}+1\big)^{2} &\qquad& \rem{$(a+b)^{2}$ avec $a=\sqrt{6}$ et $b=1$}\\
C &= \big(\sqrt{6}\big)^{2}+2\times\sqrt{6}\times 1+1^{2} &&\\
C &= 6+2\sqrt{6}+1 &&\\
C &= 7+2\sqrt{6} &&
\end{calculs}
\cas{d)}
\begin{calculs}
D &= \big(5-\sqrt{2}\big)\big(5+\sqrt{2}\big) &\qquad& \rem{$(a-b)(a+b)$ avec $a=5$ et $b=\sqrt{2}$}\\
D &= 5^{2}-\big(\sqrt{2}\big)^{2} &&\\
D &= 25-2 &&\\
D &= 23 &\qquad& \rem{les racines ont disparu}
\end{calculs}
\end{correction}

\etape{À vous de jouer}
Développer et réduire :
\[
  E=\sqrt{3}\big(2-\sqrt{3}\big) \qquad F=\big(1+\sqrt{5}\big)\big(3-\sqrt{5}\big) \qquad
  G=\big(\sqrt{7}-2\big)^{2} \qquad H=\big(4+\sqrt{3}\big)\big(4-\sqrt{3}\big)
\]
\reponse{$E=2\sqrt{3}-3$ ;\; $F=-2+2\sqrt{5}$ ;\; $G=11-4\sqrt{7}$ ;\; $H=13$.}
\end{methodebox}

\begin{methodebox}[Méthode 13 --- Écrire un quotient sans racine au dénominateur]
\etapeexemples
On multiplie le numérateur et le dénominateur par la racine du dénominateur, puisque
$\sqrt{b}\times\sqrt{b}=b$. Écrire sans racine au dénominateur :
\begin{questions}[label=\alph*)]
  \item $A=\dfrac{6}{\sqrt{3}}$
  \item $B=\dfrac{5}{2\sqrt{5}}$
\end{questions}
\begin{correction}
\cas{a)}
\begin{calculs}
A &= \dfrac{6}{\sqrt{3}} &&\\
A &= \dfrac{6\fois{\sqrt{3}}}{\sqrt{3}\fois{\sqrt{3}}} &\qquad& \rem{on multiplie en haut et en bas par $\sqrt{3}$}\\
A &= \dfrac{6\sqrt{3}}{3} &\qquad& \rem{$\sqrt{3}\times\sqrt{3}=3$}\\
A &= 2\sqrt{3} &\qquad& \rem{$6\div 3=2$}
\end{calculs}
\cas{b)}
\begin{calculs}
B &= \dfrac{5}{2\sqrt{5}} &&\\
B &= \dfrac{5\fois{\sqrt{5}}}{2\sqrt{5}\fois{\sqrt{5}}} &\qquad& \rem{seule la racine passe au numérateur}\\
B &= \dfrac{\barre{5}\sqrt{5}}{2\times\barre{5}} &\qquad& \rem{$\sqrt{5}\times\sqrt{5}=5$ ; on simplifie par $5$}\\
B &= \dfrac{\sqrt{5}}{2} &&
\end{calculs}
\end{correction}

\etape{À vous de jouer}
Écrire sans racine au dénominateur :
\[
  C=\dfrac{10}{\sqrt{5}} \qquad D=\dfrac{3}{4\sqrt{3}}
\]
\reponse{$C=2\sqrt{5}$ ;\; $D=\dfrac{\sqrt{3}}{4}$.}
\end{methodebox}

\afaire{exercices 36 à 38}

\end{document}
